% =====================================================================
%  minggu-10.tex  --  Slide Pertemuan Minggu ke-10
%  Pengantar Kriptografi (MPFA014) -- Prodi Matematika FMIPA UBB
%
%  Topik silabus : AES dan cipher blok lain
%  Sumber materi : buku_ajar/chapters/08 (Bab 8) dan chapters/09 (Bab 9)
%  Sub-CPMK      : 1.3
%
%  Kompilasi dari dalam folder ppt/:
%      pdflatex -interaction=nonstopmode minggu-10.tex   (dua kali)
% =====================================================================

\documentclass[aspectratio=169,11pt]{beamer}
% =====================================================================
%  beamer-preamble.tex  --  prasarana dan gaya slide kuliah
%
%  Mata kuliah : Pengantar Kriptografi (MPFA014)
%  Program     : Program Studi Matematika, FMIPA, Universitas Bale Bandung
%
%  Berkas ini TIDAK memuat \documentclass. Setiap berkas minggu memuatnya
%  sendiri lalu menarik berkas ini:
%
%      \documentclass[aspectratio=169,11pt]{beamer}
%      % =====================================================================
%  beamer-preamble.tex  --  prasarana dan gaya slide kuliah
%
%  Mata kuliah : Pengantar Kriptografi (MPFA014)
%  Program     : Program Studi Matematika, FMIPA, Universitas Bale Bandung
%
%  Berkas ini TIDAK memuat \documentclass. Setiap berkas minggu memuatnya
%  sendiri lalu menarik berkas ini:
%
%      \documentclass[aspectratio=169,11pt]{beamer}
%      % =====================================================================
%  beamer-preamble.tex  --  prasarana dan gaya slide kuliah
%
%  Mata kuliah : Pengantar Kriptografi (MPFA014)
%  Program     : Program Studi Matematika, FMIPA, Universitas Bale Bandung
%
%  Berkas ini TIDAK memuat \documentclass. Setiap berkas minggu memuatnya
%  sendiri lalu menarik berkas ini:
%
%      \documentclass[aspectratio=169,11pt]{beamer}
%      \input{beamer-preamble}
%      \title{...} \subtitle{...} \setWeekInfo{Minggu n}{Topik}
%
%  Berkas minggu hanya berisi \setWeekInfo dan lingkungan frame.
%  Kompilasi dijalankan DARI DALAM folder ppt/ (lihat compile-ppt.bat),
%  karena gambar buku dirujuk dengan jalur relatif ../buku_ajar/.
% =====================================================================

% ---------------------------------------------------------------- bahasa
\usepackage[utf8]{inputenc}
\usepackage[T1]{fontenc}
\usepackage[indonesian]{babel}
\usepackage{lmodern}
\usepackage{microtype}

% ------------------------------------------------------------- matematika
\usepackage{amsmath,amssymb,mathtools}

% ----------------------------------------------------------------- tabel
\usepackage{booktabs}
\usepackage{array}
\usepackage{tabularx}

% --------------------------------------------------------------- grafik
\usepackage{tikz}
\usetikzlibrary{arrows.meta,positioning,shapes.geometric,calc,fit,backgrounds}

% ---------------------------------------------------------- kode program
\usepackage{listings}
\lstdefinestyle{slidecode}{
    basicstyle=\ttfamily\scriptsize,
    keywordstyle=\color{ubbblue}\bfseries,
    commentstyle=\color{ubbgreen}\itshape,
    stringstyle=\color{ubbred},
    numbers=left,
    numberstyle=\tiny\color{gray},
    numbersep=6pt,
    frame=single,
    rulecolor=\color{ubbblue!40},
    backgroundcolor=\color{ubbblue!3},
    breaklines=true,
    showstringspaces=false,
    tabsize=4,
    columns=fullflexible,
    keepspaces=true,
}
\lstset{style=slidecode}

% =====================================================================
%  WARNA INSTITUSI
% =====================================================================
\definecolor{ubbblue}{RGB}{12, 62, 120}      % warna utama
\definecolor{ubbteal}{RGB}{0, 118, 110}      % aksen kedua
\definecolor{ubbgreen}{RGB}{22, 122, 74}     % untuk blok contoh
\definecolor{ubbred}{RGB}{168, 42, 42}       % untuk blok peringatan
\definecolor{ubbgray}{RGB}{70, 76, 84}       % teks sekunder

% =====================================================================
%  IDENTITAS MATA KULIAH
% =====================================================================
\newcommand{\CourseName}{Pengantar Kriptografi}
\newcommand{\CourseCode}{MPFA014}
\newcommand{\CourseSemester}{Semester Ganjil}
\newcommand{\CourseProgram}{Program Studi Matematika}
\newcommand{\CourseFaculty}{Fakultas Matematika dan Ilmu Pengetahuan Alam}
\newcommand{\CourseInstitution}{Universitas Bale Bandung}
\newcommand{\CourseLecturer}{Cecep Suwanda, S.Si., M.Kom.}

% Label minggu, disetel tiap berkas lewat \setWeekInfo
\newcommand{\WeekLabel}{}
\newcommand{\WeekTopic}{}
\newcommand{\setWeekInfo}[2]{%
    \renewcommand{\WeekLabel}{#1}%
    \renewcommand{\WeekTopic}{#2}%
}

% =====================================================================
%  TEMA: header, footer, warna blok
% =====================================================================
\usetheme{default}
\setbeamertemplate{navigation symbols}{}
\setbeamertemplate{blocks}[rounded][shadow=false]

\setbeamercolor{palette primary}{bg=ubbblue, fg=white}
\setbeamercolor{palette secondary}{bg=ubbgray, fg=white}
\setbeamercolor{palette tertiary}{bg=ubbteal, fg=white}

\setbeamercolor{structure}{fg=ubbblue}
\setbeamercolor{normal text}{fg=black!92}
\setbeamercolor{alerted text}{fg=ubbred}

\setbeamercolor{block title}{bg=ubbblue, fg=white}
\setbeamercolor{block body}{bg=ubbblue!6, fg=black!92}
\setbeamercolor{block title alerted}{bg=ubbred, fg=white}
\setbeamercolor{block body alerted}{bg=ubbred!6, fg=black!92}
\setbeamercolor{block title example}{bg=ubbgreen, fg=white}
\setbeamercolor{block body example}{bg=ubbgreen!6, fg=black!92}

\setbeamercolor{itemize item}{fg=ubbblue}
\setbeamercolor{itemize subitem}{fg=ubbteal}
\setbeamercolor{enumerate item}{fg=ubbblue}

\setbeamerfont{headline}{size=\scriptsize}
\setbeamerfont{footline}{size=\scriptsize}
\setbeamerfont{block title}{size=\normalsize}

% Badan blok rata kiri (bukan rata kanan-kiri): teks slide lebih mudah dibaca
% dan baris terakhir yang pendek tidak lagi memicu "Underfull \hbox".
\addtobeamertemplate{block body begin}{\raggedright}{}
\addtobeamertemplate{block alerted body begin}{\raggedright}{}
\addtobeamertemplate{block example body begin}{\raggedright}{}

% Header: nama mata kuliah dan kode MK
\setbeamertemplate{headline}{%
    \begin{beamercolorbox}[wd=\paperwidth,ht=2.6ex,dp=1.2ex,leftskip=1em,rightskip=1em]{palette primary}%
        \usebeamerfont{headline}\CourseName\ \textperiodcentered\ \CourseCode%
        \hfill \CourseProgram%
    \end{beamercolorbox}%
}

% Footer: minggu dan nomor halaman
\setbeamertemplate{footline}{%
    \begin{beamercolorbox}[wd=\paperwidth,ht=2.4ex,dp=1.2ex,leftskip=1em,rightskip=1em]{palette secondary}%
        \usebeamerfont{footline}\WeekLabel\ \textperiodcentered\ \WeekTopic%
        \hfill \insertframenumber\,/\,\inserttotalframenumber%
    \end{beamercolorbox}%
}

% Halaman judul: judul, subjudul, lalu blok identitas mata kuliah
\setbeamertemplate{title page}{%
    \begin{centering}
        \vskip1.0em
        {\usebeamerfont{title}\usebeamercolor[fg]{title}\inserttitle\par}
        \vskip0.3em
        {\usebeamerfont{subtitle}\usebeamercolor[fg]{subtitle}\insertsubtitle\par}
        \vskip1.2em
        \textcolor{ubbblue}{\rule{0.72\textwidth}{0.8pt}}
        \vskip1.2em
        {\small
            \textbf{\CourseName}\ (\CourseCode)\par
            \vskip0.4em
            \CourseProgram\par
            \CourseFaculty\par
            \CourseInstitution\par
            \vskip0.6em
            \CourseSemester\ \textperiodcentered\ \WeekLabel\par
            \vskip0.8em
            Dosen Pengampu: \CourseLecturer\par
        }
        \vskip1.0em
        \textcolor{ubbblue}{\rule{0.72\textwidth}{0.4pt}}
    \end{centering}
}

% =====================================================================
%  PERINTAH BANTU
% =====================================================================
\newcommand{\R}{\mathbb{R}}
\newcommand{\N}{\mathbb{N}}
\newcommand{\Z}{\mathbb{Z}}
\newcommand{\code}[1]{\texttt{#1}}

% Rujukan gambar dari buku: \bookfig[lebar]{NN}{namaberkas.png}
% Jalur relatif terhadap folder ppt/ (dikompilasi dari sana).
\newcommand{\bookfig}[3][0.80\linewidth]{%
    \includegraphics[width=#1]{../buku_ajar/chapters/#2/figures/#3}%
}

% Keterangan sumber gambar, sepadan dengan \sumber di buku
\newcommand{\figsumber}[1]{%
    \par\vspace{1pt}{\tiny\itshape Sumber: #1\par}%
}

% Penanda blok materi yang sering dipakai di slide
\newcommand{\slidejudul}[1]{{\large\bfseries\textcolor{ubbblue}{#1}}\par\vspace{2pt}}
\newcommand{\kunci}[1]{\textcolor{ubbred}{\textbf{#1}}}

%      \title{...} \subtitle{...} \setWeekInfo{Minggu n}{Topik}
%
%  Berkas minggu hanya berisi \setWeekInfo dan lingkungan frame.
%  Kompilasi dijalankan DARI DALAM folder ppt/ (lihat compile-ppt.bat),
%  karena gambar buku dirujuk dengan jalur relatif ../buku_ajar/.
% =====================================================================

% ---------------------------------------------------------------- bahasa
\usepackage[utf8]{inputenc}
\usepackage[T1]{fontenc}
\usepackage[indonesian]{babel}
\usepackage{lmodern}
\usepackage{microtype}

% ------------------------------------------------------------- matematika
\usepackage{amsmath,amssymb,mathtools}

% ----------------------------------------------------------------- tabel
\usepackage{booktabs}
\usepackage{array}
\usepackage{tabularx}

% --------------------------------------------------------------- grafik
\usepackage{tikz}
\usetikzlibrary{arrows.meta,positioning,shapes.geometric,calc,fit,backgrounds}

% ---------------------------------------------------------- kode program
\usepackage{listings}
\lstdefinestyle{slidecode}{
    basicstyle=\ttfamily\scriptsize,
    keywordstyle=\color{ubbblue}\bfseries,
    commentstyle=\color{ubbgreen}\itshape,
    stringstyle=\color{ubbred},
    numbers=left,
    numberstyle=\tiny\color{gray},
    numbersep=6pt,
    frame=single,
    rulecolor=\color{ubbblue!40},
    backgroundcolor=\color{ubbblue!3},
    breaklines=true,
    showstringspaces=false,
    tabsize=4,
    columns=fullflexible,
    keepspaces=true,
}
\lstset{style=slidecode}

% =====================================================================
%  WARNA INSTITUSI
% =====================================================================
\definecolor{ubbblue}{RGB}{12, 62, 120}      % warna utama
\definecolor{ubbteal}{RGB}{0, 118, 110}      % aksen kedua
\definecolor{ubbgreen}{RGB}{22, 122, 74}     % untuk blok contoh
\definecolor{ubbred}{RGB}{168, 42, 42}       % untuk blok peringatan
\definecolor{ubbgray}{RGB}{70, 76, 84}       % teks sekunder

% =====================================================================
%  IDENTITAS MATA KULIAH
% =====================================================================
\newcommand{\CourseName}{Pengantar Kriptografi}
\newcommand{\CourseCode}{MPFA014}
\newcommand{\CourseSemester}{Semester Ganjil}
\newcommand{\CourseProgram}{Program Studi Matematika}
\newcommand{\CourseFaculty}{Fakultas Matematika dan Ilmu Pengetahuan Alam}
\newcommand{\CourseInstitution}{Universitas Bale Bandung}
\newcommand{\CourseLecturer}{Cecep Suwanda, S.Si., M.Kom.}

% Label minggu, disetel tiap berkas lewat \setWeekInfo
\newcommand{\WeekLabel}{}
\newcommand{\WeekTopic}{}
\newcommand{\setWeekInfo}[2]{%
    \renewcommand{\WeekLabel}{#1}%
    \renewcommand{\WeekTopic}{#2}%
}

% =====================================================================
%  TEMA: header, footer, warna blok
% =====================================================================
\usetheme{default}
\setbeamertemplate{navigation symbols}{}
\setbeamertemplate{blocks}[rounded][shadow=false]

\setbeamercolor{palette primary}{bg=ubbblue, fg=white}
\setbeamercolor{palette secondary}{bg=ubbgray, fg=white}
\setbeamercolor{palette tertiary}{bg=ubbteal, fg=white}

\setbeamercolor{structure}{fg=ubbblue}
\setbeamercolor{normal text}{fg=black!92}
\setbeamercolor{alerted text}{fg=ubbred}

\setbeamercolor{block title}{bg=ubbblue, fg=white}
\setbeamercolor{block body}{bg=ubbblue!6, fg=black!92}
\setbeamercolor{block title alerted}{bg=ubbred, fg=white}
\setbeamercolor{block body alerted}{bg=ubbred!6, fg=black!92}
\setbeamercolor{block title example}{bg=ubbgreen, fg=white}
\setbeamercolor{block body example}{bg=ubbgreen!6, fg=black!92}

\setbeamercolor{itemize item}{fg=ubbblue}
\setbeamercolor{itemize subitem}{fg=ubbteal}
\setbeamercolor{enumerate item}{fg=ubbblue}

\setbeamerfont{headline}{size=\scriptsize}
\setbeamerfont{footline}{size=\scriptsize}
\setbeamerfont{block title}{size=\normalsize}

% Badan blok rata kiri (bukan rata kanan-kiri): teks slide lebih mudah dibaca
% dan baris terakhir yang pendek tidak lagi memicu "Underfull \hbox".
\addtobeamertemplate{block body begin}{\raggedright}{}
\addtobeamertemplate{block alerted body begin}{\raggedright}{}
\addtobeamertemplate{block example body begin}{\raggedright}{}

% Header: nama mata kuliah dan kode MK
\setbeamertemplate{headline}{%
    \begin{beamercolorbox}[wd=\paperwidth,ht=2.6ex,dp=1.2ex,leftskip=1em,rightskip=1em]{palette primary}%
        \usebeamerfont{headline}\CourseName\ \textperiodcentered\ \CourseCode%
        \hfill \CourseProgram%
    \end{beamercolorbox}%
}

% Footer: minggu dan nomor halaman
\setbeamertemplate{footline}{%
    \begin{beamercolorbox}[wd=\paperwidth,ht=2.4ex,dp=1.2ex,leftskip=1em,rightskip=1em]{palette secondary}%
        \usebeamerfont{footline}\WeekLabel\ \textperiodcentered\ \WeekTopic%
        \hfill \insertframenumber\,/\,\inserttotalframenumber%
    \end{beamercolorbox}%
}

% Halaman judul: judul, subjudul, lalu blok identitas mata kuliah
\setbeamertemplate{title page}{%
    \begin{centering}
        \vskip1.0em
        {\usebeamerfont{title}\usebeamercolor[fg]{title}\inserttitle\par}
        \vskip0.3em
        {\usebeamerfont{subtitle}\usebeamercolor[fg]{subtitle}\insertsubtitle\par}
        \vskip1.2em
        \textcolor{ubbblue}{\rule{0.72\textwidth}{0.8pt}}
        \vskip1.2em
        {\small
            \textbf{\CourseName}\ (\CourseCode)\par
            \vskip0.4em
            \CourseProgram\par
            \CourseFaculty\par
            \CourseInstitution\par
            \vskip0.6em
            \CourseSemester\ \textperiodcentered\ \WeekLabel\par
            \vskip0.8em
            Dosen Pengampu: \CourseLecturer\par
        }
        \vskip1.0em
        \textcolor{ubbblue}{\rule{0.72\textwidth}{0.4pt}}
    \end{centering}
}

% =====================================================================
%  PERINTAH BANTU
% =====================================================================
\newcommand{\R}{\mathbb{R}}
\newcommand{\N}{\mathbb{N}}
\newcommand{\Z}{\mathbb{Z}}
\newcommand{\code}[1]{\texttt{#1}}

% Rujukan gambar dari buku: \bookfig[lebar]{NN}{namaberkas.png}
% Jalur relatif terhadap folder ppt/ (dikompilasi dari sana).
\newcommand{\bookfig}[3][0.80\linewidth]{%
    \includegraphics[width=#1]{../buku_ajar/chapters/#2/figures/#3}%
}

% Keterangan sumber gambar, sepadan dengan \sumber di buku
\newcommand{\figsumber}[1]{%
    \par\vspace{1pt}{\tiny\itshape Sumber: #1\par}%
}

% Penanda blok materi yang sering dipakai di slide
\newcommand{\slidejudul}[1]{{\large\bfseries\textcolor{ubbblue}{#1}}\par\vspace{2pt}}
\newcommand{\kunci}[1]{\textcolor{ubbred}{\textbf{#1}}}

%      \title{...} \subtitle{...} \setWeekInfo{Minggu n}{Topik}
%
%  Berkas minggu hanya berisi \setWeekInfo dan lingkungan frame.
%  Kompilasi dijalankan DARI DALAM folder ppt/ (lihat compile-ppt.bat),
%  karena gambar buku dirujuk dengan jalur relatif ../buku_ajar/.
% =====================================================================

% ---------------------------------------------------------------- bahasa
\usepackage[utf8]{inputenc}
\usepackage[T1]{fontenc}
\usepackage[indonesian]{babel}
\usepackage{lmodern}
\usepackage{microtype}

% ------------------------------------------------------------- matematika
\usepackage{amsmath,amssymb,mathtools}

% ----------------------------------------------------------------- tabel
\usepackage{booktabs}
\usepackage{array}
\usepackage{tabularx}

% --------------------------------------------------------------- grafik
\usepackage{tikz}
\usetikzlibrary{arrows.meta,positioning,shapes.geometric,calc,fit,backgrounds}

% ---------------------------------------------------------- kode program
\usepackage{listings}
\lstdefinestyle{slidecode}{
    basicstyle=\ttfamily\scriptsize,
    keywordstyle=\color{ubbblue}\bfseries,
    commentstyle=\color{ubbgreen}\itshape,
    stringstyle=\color{ubbred},
    numbers=left,
    numberstyle=\tiny\color{gray},
    numbersep=6pt,
    frame=single,
    rulecolor=\color{ubbblue!40},
    backgroundcolor=\color{ubbblue!3},
    breaklines=true,
    showstringspaces=false,
    tabsize=4,
    columns=fullflexible,
    keepspaces=true,
}
\lstset{style=slidecode}

% =====================================================================
%  WARNA INSTITUSI
% =====================================================================
\definecolor{ubbblue}{RGB}{12, 62, 120}      % warna utama
\definecolor{ubbteal}{RGB}{0, 118, 110}      % aksen kedua
\definecolor{ubbgreen}{RGB}{22, 122, 74}     % untuk blok contoh
\definecolor{ubbred}{RGB}{168, 42, 42}       % untuk blok peringatan
\definecolor{ubbgray}{RGB}{70, 76, 84}       % teks sekunder

% =====================================================================
%  IDENTITAS MATA KULIAH
% =====================================================================
\newcommand{\CourseName}{Pengantar Kriptografi}
\newcommand{\CourseCode}{MPFA014}
\newcommand{\CourseSemester}{Semester Ganjil}
\newcommand{\CourseProgram}{Program Studi Matematika}
\newcommand{\CourseFaculty}{Fakultas Matematika dan Ilmu Pengetahuan Alam}
\newcommand{\CourseInstitution}{Universitas Bale Bandung}
\newcommand{\CourseLecturer}{Cecep Suwanda, S.Si., M.Kom.}

% Label minggu, disetel tiap berkas lewat \setWeekInfo
\newcommand{\WeekLabel}{}
\newcommand{\WeekTopic}{}
\newcommand{\setWeekInfo}[2]{%
    \renewcommand{\WeekLabel}{#1}%
    \renewcommand{\WeekTopic}{#2}%
}

% =====================================================================
%  TEMA: header, footer, warna blok
% =====================================================================
\usetheme{default}
\setbeamertemplate{navigation symbols}{}
\setbeamertemplate{blocks}[rounded][shadow=false]

\setbeamercolor{palette primary}{bg=ubbblue, fg=white}
\setbeamercolor{palette secondary}{bg=ubbgray, fg=white}
\setbeamercolor{palette tertiary}{bg=ubbteal, fg=white}

\setbeamercolor{structure}{fg=ubbblue}
\setbeamercolor{normal text}{fg=black!92}
\setbeamercolor{alerted text}{fg=ubbred}

\setbeamercolor{block title}{bg=ubbblue, fg=white}
\setbeamercolor{block body}{bg=ubbblue!6, fg=black!92}
\setbeamercolor{block title alerted}{bg=ubbred, fg=white}
\setbeamercolor{block body alerted}{bg=ubbred!6, fg=black!92}
\setbeamercolor{block title example}{bg=ubbgreen, fg=white}
\setbeamercolor{block body example}{bg=ubbgreen!6, fg=black!92}

\setbeamercolor{itemize item}{fg=ubbblue}
\setbeamercolor{itemize subitem}{fg=ubbteal}
\setbeamercolor{enumerate item}{fg=ubbblue}

\setbeamerfont{headline}{size=\scriptsize}
\setbeamerfont{footline}{size=\scriptsize}
\setbeamerfont{block title}{size=\normalsize}

% Badan blok rata kiri (bukan rata kanan-kiri): teks slide lebih mudah dibaca
% dan baris terakhir yang pendek tidak lagi memicu "Underfull \hbox".
\addtobeamertemplate{block body begin}{\raggedright}{}
\addtobeamertemplate{block alerted body begin}{\raggedright}{}
\addtobeamertemplate{block example body begin}{\raggedright}{}

% Header: nama mata kuliah dan kode MK
\setbeamertemplate{headline}{%
    \begin{beamercolorbox}[wd=\paperwidth,ht=2.6ex,dp=1.2ex,leftskip=1em,rightskip=1em]{palette primary}%
        \usebeamerfont{headline}\CourseName\ \textperiodcentered\ \CourseCode%
        \hfill \CourseProgram%
    \end{beamercolorbox}%
}

% Footer: minggu dan nomor halaman
\setbeamertemplate{footline}{%
    \begin{beamercolorbox}[wd=\paperwidth,ht=2.4ex,dp=1.2ex,leftskip=1em,rightskip=1em]{palette secondary}%
        \usebeamerfont{footline}\WeekLabel\ \textperiodcentered\ \WeekTopic%
        \hfill \insertframenumber\,/\,\inserttotalframenumber%
    \end{beamercolorbox}%
}

% Halaman judul: judul, subjudul, lalu blok identitas mata kuliah
\setbeamertemplate{title page}{%
    \begin{centering}
        \vskip1.0em
        {\usebeamerfont{title}\usebeamercolor[fg]{title}\inserttitle\par}
        \vskip0.3em
        {\usebeamerfont{subtitle}\usebeamercolor[fg]{subtitle}\insertsubtitle\par}
        \vskip1.2em
        \textcolor{ubbblue}{\rule{0.72\textwidth}{0.8pt}}
        \vskip1.2em
        {\small
            \textbf{\CourseName}\ (\CourseCode)\par
            \vskip0.4em
            \CourseProgram\par
            \CourseFaculty\par
            \CourseInstitution\par
            \vskip0.6em
            \CourseSemester\ \textperiodcentered\ \WeekLabel\par
            \vskip0.8em
            Dosen Pengampu: \CourseLecturer\par
        }
        \vskip1.0em
        \textcolor{ubbblue}{\rule{0.72\textwidth}{0.4pt}}
    \end{centering}
}

% =====================================================================
%  PERINTAH BANTU
% =====================================================================
\newcommand{\R}{\mathbb{R}}
\newcommand{\N}{\mathbb{N}}
\newcommand{\Z}{\mathbb{Z}}
\newcommand{\code}[1]{\texttt{#1}}

% Rujukan gambar dari buku: \bookfig[lebar]{NN}{namaberkas.png}
% Jalur relatif terhadap folder ppt/ (dikompilasi dari sana).
\newcommand{\bookfig}[3][0.80\linewidth]{%
    \includegraphics[width=#1]{../buku_ajar/chapters/#2/figures/#3}%
}

% Keterangan sumber gambar, sepadan dengan \sumber di buku
\newcommand{\figsumber}[1]{%
    \par\vspace{1pt}{\tiny\itshape Sumber: #1\par}%
}

% Penanda blok materi yang sering dipakai di slide
\newcommand{\slidejudul}[1]{{\large\bfseries\textcolor{ubbblue}{#1}}\par\vspace{2pt}}
\newcommand{\kunci}[1]{\textcolor{ubbred}{\textbf{#1}}}


\title{AES dan Cipher Blok Lain}
\subtitle{Rijndael, state $4 \times 4$, dan tujuh cipher blok pembanding}
\setWeekInfo{Minggu 10}{AES dan cipher blok lain}

\begin{document}

% ----------------------------------------------------------- halaman judul
\begin{frame}[plain]
    \titlepage
\end{frame}

% ------------------------------------------------------------------ peta
\begin{frame}{Peta Materi Pertemuan Ini}
    \begin{columns}[T,onlytextwidth]
        \begin{column}{0.56\textwidth}
            \slidejudul{Empat bagian}
            \begin{enumerate}
                \item Dari DES ke AES: syarat, finalis, dan varian.
                \item State $4 \times 4$ dan empat transformasi putaran.
                \item Contoh kerja satu state menembus tiga transformasi.
                \item Cipher blok lain: GOST, RC5, RC6, Blowfish,
                      Serpent, Twofish, dan MARS.
            \end{enumerate}
        \end{column}
        \begin{column}{0.40\textwidth}
            \begin{block}{Referensi}
                \footnotesize
                Bab 8 dan Bab 9 buku ajar; Munir, \textit{IF4020
                Kriptografi}, ITB, 2026 (modul 10 dan 11)
            \end{block}
            \begin{exampleblock}{Dua bab, satu pertemuan}
                \footnotesize
                Bab 8 membahas AES secara utuh; Bab 9 menjadi kamus
                pembanding tujuh cipher blok lain.
            \end{exampleblock}
        \end{column}
    \end{columns}
\end{frame}

% ---------------------------------------------------- capaian pembelajaran
\begin{frame}{Capaian Pembelajaran}
    \begin{block}{Sub-CPMK 1.3}
        \footnotesize
        Mengenal block cipher dan mode operasinya (ECB, CBC, CFB, OFB,
        CTR), jaringan Feistel, DES, Triple DES, AES, serta beberapa block
        cipher lain (GOST, RC5, RC6, Blowfish, Serpent, Twofish, dan MARS).
    \end{block}
    \vspace{0.15cm}
    \begin{exampleblock}{Setelah pertemuan ini, mahasiswa mampu}
        \footnotesize
        \begin{itemize}
            \item menjelaskan alasan lahirnya AES dan persyaratan calon
                  pemenang sayembara;
            \item menjelaskan peran state $4 \times 4$ dan keempat
                  transformasi putaran AES;
            \item mengidentifikasi ukuran blok, ukuran kunci, dan ciri
                  pembeda tiap cipher blok lain;
            \item membandingkan AES dengan GOST, RC5, RC6, Blowfish,
                  Serpent, Twofish, dan MARS.
        \end{itemize}
    \end{exampleblock}
\end{frame}

% =====================================================================
\section{Dari DES ke AES}
% =====================================================================

\begin{frame}{Mengapa Lahir Standar Baru}
    \footnotesize
    \begin{itemize}
        \item Kunci DES yang hanya 56 bit dapat dihabiskan dengan
              penelusuran menyeluruh.
        \item Diperlukan cipher blok baru yang tahan terhadap serangan
              semacam itu.
        \item Sayembara terbuka: setiap rancangan dinilai di muka umum,
              tanpa bagian yang dirahasiakan.
        \item Oktober 2000: Rijndael karya Joan Daemen dan Vincent Rijmen
              dipilih sebagai pemenang.
        \item November 2001: Rijndael ditetapkan sebagai \textit{Advanced
              Encryption Standard} (AES).
    \end{itemize}
    \vspace{0.1cm}
    \begin{alertblock}{Nama dan asal}
        \footnotesize
        ``Rijndael'' dibaca \textit{rhine-doll}; AES adalah nama standar
        yang melekat pada rancangan itu.
    \end{alertblock}
\end{frame}

\begin{frame}{Lima Persyaratan Calon AES}
    \footnotesize
    \begin{enumerate}
        \item Termasuk kriptografi simetri berbasis cipher blok.
        \item Seluruh rancangan publik; tidak ada bagian yang
              dirahasiakan.
        \item Panjang kunci fleksibel: 128, 192, dan 256 bit.
        \item Ukuran blok yang dienkripsi 128 bit.
        \item Dapat diimplementasikan sebagai perangkat lunak maupun
              perangkat keras.
    \end{enumerate}
    \vspace{0.2cm}
    \begin{alertblock}{Konsekuensi syarat kedua}
        \footnotesize
        Keamanan tidak boleh bergantung pada kerahasiaan algoritma.
        Yang dirahasiakan hanyalah kuncinya.
    \end{alertblock}
\end{frame}

\begin{frame}{Lima Finalis dan Pemenang}
    \footnotesize
    \begin{columns}[T,onlytextwidth]
        \begin{column}{0.64\textwidth}
            \begin{tabularx}{\linewidth}{@{}l >{\raggedright\arraybackslash}X c@{}}
                \toprule
                \textbf{Algoritma} & \textbf{Perancang dan asal} & \textbf{Suara} \\
                \midrule
                Rijndael & Daemen dan Rijmen (Belgia) & 86 \\
                Serpent & Anderson, Biham, Knudsen & 59 \\
                Twofish & Tim yang diketuai Schneier & 31 \\
                RC6 & RSA Laboratories (Amerika) & 23 \\
                MARS & IBM & 13 \\
                \bottomrule
            \end{tabularx}
            \vspace{0.15cm}
            Rijndael memperoleh 86 suara, terbanyak di antara kelima
            finalis, lalu ditetapkan sebagai AES.
        \end{column}
        \begin{column}{0.32\textwidth}
            \centering
            \bookfig[0.95\linewidth]{08}{perancang-aes.png}
            \figsumber{Munir, \textit{IF4020 Kriptografi}, ITB.}
        \end{column}
    \end{columns}
\end{frame}

\begin{frame}{Tiga Varian AES}
    \footnotesize
    Satuan panjang kunci dan blok adalah \textit{word} 32 bit. Parameter
    varian ditulis $N_k$ (word kunci), $N_b$ (word blok), dan $N_r$
    (jumlah putaran).
    \vspace{0.1cm}
    \begin{center}
    \begin{tabular}{@{}l c c c c@{}}
        \toprule
        \textbf{Varian} & \textbf{Panjang kunci} & $N_k$ & $N_b$ & $N_r$ \\
        \midrule
        AES-128 & 128 bit & 4 word & 4 word & 10 \\
        AES-192 & 192 bit & 6 word & 4 word & 12 \\
        AES-256 & 256 bit & 8 word & 4 word & 14 \\
        \bottomrule
    \end{tabular}
    \end{center}
    \vspace{0.1cm}
    \begin{itemize}
        \item Rijndael membolehkan kunci 128 sampai 256 bit dengan
              kenaikan 32 bit.
        \item Panjang kunci dan ukuran blok pada Rijndael dapat dipilih
              terpisah; AES mengunci blok pada 128 bit.
        \item Dalam praktik, AES-128 dan AES-256 yang lazim dipakai;
              kunci 192 bit jarang.
    \end{itemize}
\end{frame}

\begin{frame}{Ruang Kunci dan Penelusuran Menyeluruh}
    \begin{block}{Banyaknya kemungkinan kunci AES-128}
        \[
            2^{128} = 3{,}4 \times 10^{38}
        \]
    \end{block}
    \vspace{0.1cm}
    \begin{columns}[T,onlytextwidth]
        \begin{column}{0.47\textwidth}
            \begin{exampleblock}{Satu juta kunci per detik}
                \footnotesize
                Waktu rata-rata mencoba seluruh kunci:
                \[
                    5{,}4 \times 10^{24} \text{ tahun}.
                \]
            \end{exampleblock}
        \end{column}
        \begin{column}{0.47\textwidth}
            \begin{exampleblock}{Satu miliar kunci per detik}
                \footnotesize
                Waktu rata-rata turun menjadi
                \[
                    5{,}4 \times 10^{21} \text{ tahun}.
                \]
            \end{exampleblock}
        \end{column}
    \end{columns}
    \vspace{0.1cm}
    \begin{alertblock}{Kesimpulan}
        \footnotesize
        Bahkan dengan laju yang jauh melampaui komputer biasa, penelusuran
        kunci menyeluruh tetap tidak praktis.
    \end{alertblock}
\end{frame}

\begin{frame}{Ciri Rancangan AES}
    \footnotesize
    \begin{columns}[T,onlytextwidth]
        \begin{column}{0.48\textwidth}
            \begin{block}{Bekerja pada byte}
                Seluruh operasi memandang data sebagai satuan 8 bit, bukan
                rangkaian bit tunggal.
            \end{block}
            \begin{block}{Bukan jaringan Feistel}
                Setiap putaran mengolah seluruh state, bukan hanya separuh
                bagian seperti pada jaringan Feistel.
            \end{block}
        \end{column}
        \begin{column}{0.48\textwidth}
            \begin{block}{Kunci putaran berbeda}
                Setiap putaran memakai kunci internal tersendiri yang
                diturunkan dari kunci pengguna.
            \end{block}
            \begin{block}{Substitusi dan permutasi}
                Tiap putaran menggabungkan operasi substitusi dengan
                operasi permutasi.
            \end{block}
        \end{column}
    \end{columns}
    \vspace{0.15cm}
    Perbedaan cara bekerja inilah yang membuat AES tidak dapat dipandang
    sebagai DES dengan kunci yang diperpanjang.
\end{frame}

% =====================================================================
\section{State dan Empat Transformasi}
% =====================================================================

\begin{frame}{State: Matriks $4 \times 4$}
    \footnotesize
    Selama enkripsi, data disimpan dalam \texttt{state}, yaitu matriks
    byte berukuran $N_{\text{rows}} \times N_{\text{cols}}$. Untuk blok
    128 bit, ukurannya $4 \times 4$.
    \vspace{0.1cm}
    \begin{columns}[T,onlytextwidth]
        \begin{column}{0.46\textwidth}
            \begin{block}{Pemetaan masukan}
                Enam belas byte $in_0, \dots, in_{15}$ disalin kolom demi
                kolom:
                \[
                    s_{r,c} = in_{r + 4c}, \qquad r,c \in \{0,1,2,3\}.
                \]
            \end{block}
        \end{column}
        \begin{column}{0.46\textwidth}
            \centering
            \[
            \begin{bmatrix}
                in_0 & in_4 & in_8  & in_{12} \\
                in_1 & in_5 & in_9  & in_{13} \\
                in_2 & in_6 & in_{10} & in_{14} \\
                in_3 & in_7 & in_{11} & in_{15}
            \end{bmatrix}
            \]
        \end{column}
    \end{columns}
    \vspace{0.1cm}
    \begin{alertblock}{Kekeliruan yang sering terjadi}
        \footnotesize
        Empat byte pertama mengisi kolom pertama, bukan baris pertama.
        Urutan yang tertukar akan mengacaukan seluruh hasil.
    \end{alertblock}
\end{frame}

\begin{frame}{SubBytes}
    \footnotesize
    \begin{columns}[T,onlytextwidth]
        \begin{column}{0.58\textwidth}
            \begin{block}{Operasi}
                Setiap byte diganti nilai dari tabel substitusi
                \textbf{S-box}. Posisi byte tidak berubah; hanya nilainya
                yang berubah:
                \[
                    s'_{r,c} = \text{S-box}(x,y), \quad s_{r,c} = xy.
                \]
            \end{block}
            \begin{exampleblock}{Contoh buku}
                Byte \code{23} memiliki digit atas 2 dan digit bawah 3.
                S-box baris 2 kolom 3 bernilai \code{26}, sehingga
                $\code{23} \to \code{26}$.
            \end{exampleblock}
        \end{column}
        \begin{column}{0.38\textwidth}
            \centering
            \scriptsize
            \begin{tabular}{@{}c ccc@{}}
                \toprule
                & \textbf{2} & \textbf{3} & \textbf{4} \\
                \midrule
                \textbf{1} & C9 & 7D & FA \\
                \textbf{2} & 93 & \textbf{26} & 36 \\
                \textbf{3} & \textbf{23} & C3 & 18 \\
                \bottomrule
            \end{tabular}
        \end{column}
    \end{columns}
    \vspace{0.1cm}
    \begin{alertblock}{Peran dalam AES}
        \footnotesize
        SubBytes adalah satu-satunya transformasi nonlinear; di sinilah
        sifat \textit{confusion} dibangun.
    \end{alertblock}
\end{frame}

\begin{frame}{ShiftRows}
    \footnotesize
    \begin{columns}[T,onlytextwidth]
        \begin{column}{0.52\textwidth}
            \begin{block}{Operasi}
                Tiga baris terakhir digeser siklik ke kiri. Baris $r$
                bergeser sejauh $r$ byte:
                \[
                    s'_{r,c} = s_{r,\,(c+r) \bmod 4}.
                \]
            \end{block}
            \begin{itemize}
                \item Baris 0 tetap.
                \item Baris 1 bergeser 1 byte.
                \item Baris 2 bergeser 2 byte.
                \item Baris 3 bergeser 3 byte.
            \end{itemize}
        \end{column}
        \begin{column}{0.44\textwidth}
            \centering
            \bookfig[0.98\linewidth]{08}{shiftrows-aes.png}
            \figsumber{Munir, \textit{IF4020 Kriptografi}, ITB.}
        \end{column}
    \end{columns}
    \vspace{0.05cm}
    \begin{alertblock}{Peran dalam AES}
        \footnotesize
        Setelah pergeseran, satu kolom memuat byte dari empat kolom asal
        yang berbeda. Penyebaran antarposisi ini mendukung
        \textit{diffusion}.
    \end{alertblock}
\end{frame}

\begin{frame}{MixColumns}
    \footnotesize
    \begin{block}{Operasi}
        Setiap kolom dikalikan dengan matriks tetap. Seluruh perkalian dan
        penjumlahan berlangsung di $\text{GF}(2^8)$:
        \[
        \begin{bmatrix}
            s'_{0,c} \\ s'_{1,c} \\ s'_{2,c} \\ s'_{3,c}
        \end{bmatrix}
        =
        \begin{bmatrix}
            \code{02} & \code{03} & \code{01} & \code{01} \\
            \code{01} & \code{02} & \code{03} & \code{01} \\
            \code{01} & \code{01} & \code{02} & \code{03} \\
            \code{03} & \code{01} & \code{01} & \code{02}
        \end{bmatrix}
        \begin{bmatrix}
            s_{0,c} \\ s_{1,c} \\ s_{2,c} \\ s_{3,c}
        \end{bmatrix}.
        \]
    \end{block}
    \vspace{0.05cm}
    \begin{itemize}
        \item Penjumlahan di $\text{GF}(2^8)$ sama dengan XOR.
        \item Perkalian direduksi oleh polinom
              $x^8 + x^4 + x^3 + x + 1$.
        \item Setiap byte masukan memengaruhi keempat byte keluaran pada
              kolom yang sama.
    \end{itemize}
\end{frame}

\begin{frame}{AddRoundKey}
    \footnotesize
    \begin{block}{Operasi}
        State di-XOR dengan kunci putaran pada posisi yang sama:
        \[
            s'_{r,c} = s_{r,c} \oplus k_{r,c}.
        \]
        Tidak ada aritmetika medan di sini; operasinya murni XOR byte.
    \end{block}
    \vspace{0.1cm}
    \begin{itemize}
        \item Hanya transformasi yang menyisipkan kunci ke dalam state.
        \item Bersifat membalik dirinya sendiri: XOR dua kali dengan kunci
              yang sama mengembalikan nilai semula.
        \item Kunci putaran $k_{r,c}$ berasal dari proses
              \textit{key expansion}, bukan langsung dari kunci pengguna.
    \end{itemize}
    \vspace{0.05cm}
    \begin{alertblock}{Tanpa AddRoundKey}
        \footnotesize
        SubBytes, ShiftRows, dan MixColumns tidak mengenal kunci sama
        sekali; tanpa AddRoundKey, enkripsi tidak bergantung pada kunci.
    \end{alertblock}
\end{frame}

\begin{frame}{Alur Enkripsi AES-128}
    \footnotesize
    \begin{columns}[T,onlytextwidth]
        \begin{column}{0.50\textwidth}
            \begin{enumerate}
                \item \textit{AddRoundKey} awal: state di-XOR dengan kunci
                      pengguna.
                \item Sebanyak $N_r - 1 = 9$ putaran penuh, masing-masing
                      SubBytes, ShiftRows, MixColumns, lalu AddRoundKey.
                \item Putaran terakhir: SubBytes, ShiftRows, dan
                      AddRoundKey, \emph{tanpa} MixColumns.
            \end{enumerate}
            \vspace{0.1cm}
            \begin{alertblock}{Putaran terakhir}
                \footnotesize
                Hanya tiga transformasi yang dijalankan; dekripsi membalik
                urutan itu dengan transformasi inversnya.
            \end{alertblock}
        \end{column}
        \begin{column}{0.46\textwidth}
            \centering
            \bookfig[0.98\linewidth]{08}{alur-enkripsi-aes.png}
            \figsumber{Munir, \textit{IF4020 Kriptografi}, ITB.}
        \end{column}
    \end{columns}
\end{frame}

\begin{frame}{Key Expansion}
    \footnotesize
    Semua kunci putaran diturunkan dari satu \textit{cipher key} melalui
    fungsi \textit{KeyExpansion}. Satuannya \textit{word} 32 bit; AES-128
    menghasilkan larik $w[0]$ sampai $w[43]$, yaitu 11 kunci putaran.
    \vspace{0.1cm}
    \begin{enumerate}
        \item Salin kunci pengguna ke $w[0], w[1], w[2], w[3]$.
        \item Untuk $i = 4$ sampai 43, simpan $w[i-1]$ di peubah
              \textit{temp}.
        \item Bila $i$ kelipatan 4, jalankan fungsi $g$: RotWord
              (geser satu byte), SubWord (substitusi S-box), lalu XOR
              dengan $Rcon[i/4]$.
        \item $w[i] = w[i-4] \oplus \textit{temp}$.
    \end{enumerate}
    \vspace{0.05cm}
    \begin{alertblock}{Konstanta putaran}
        \footnotesize
        $RC[1]$ sampai $RC[10]$: \code{01}, \code{02}, \code{04},
        \code{08}, \code{10}, \code{20}, \code{40}, \code{80},
        \code{1B}, \code{36}.
    \end{alertblock}
\end{frame}

\begin{frame}{Contoh Key Expansion dari Buku}
    \footnotesize
    Kunci pengguna adalah teks ASCII ``Two One Nine Two'':
    \[
        (\code{54},\code{77},\code{6F},\code{20},\code{4F},\code{6E},
        \code{65},\code{20},\code{4E},\code{69},\code{6E},\code{65},
        \code{20},\code{54},\code{77},\code{6F}).
    \]
    Word terakhir kunci adalah $w[3] = (\code{20},\code{54},\code{77},
    \code{6F})$. Pembangkitan $w[4]$ berjalan sebagai berikut.
    \vspace{0.05cm}
    \begin{align*}
        \text{RotWord} &: (\code{20},\code{54},\code{77},\code{6F})
        \longrightarrow (\code{54},\code{77},\code{6F},\code{20}), \\
        \text{SubWord} &: (\code{54},\code{77},\code{6F},\code{20})
        \longrightarrow (\code{20},\code{F5},\code{A8},\code{B7}), \\
        \text{XOR } Rcon[1] &: (\code{20},\code{F5},\code{A8},\code{B7})
        \oplus (\code{01},\code{00},\code{00},\code{00})
        = (\code{21},\code{F5},\code{A8},\code{B7}), \\
        w[4] &= w[0] \oplus g(w[3]) \\
        &= (\code{54},\code{77},\code{6F},\code{20})
        \oplus (\code{21},\code{F5},\code{A8},\code{B7})
        = (\code{75},\code{82},\code{C7},\code{97}).
    \end{align*}
    \vspace{-0.1cm}
    \begin{exampleblock}{Lanjutan}
        \scriptsize
        $w[5] = w[4] \oplus w[1] = \code{3AECA2B7}$,\quad
        $w[6] = w[5] \oplus w[2] = \code{7485CCD2}$,\quad
        $w[7] = w[6] \oplus w[3] = \code{54D1BBBD}$.
    \end{exampleblock}
\end{frame}

\begin{frame}{Dekripsi AES}
    \footnotesize
    Dekripsi membalik urutan enkripsi dengan tiga transformasi invers.
    \vspace{0.1cm}
    \begin{columns}[T,onlytextwidth]
        \begin{column}{0.48\textwidth}
            \begin{block}{Urutan kerja}
                \begin{itemize}
                    \item Cipherteks di-XOR dengan kunci putaran terakhir.
                    \item Setiap putaran: InvShiftRows, InvSubBytes,
                          AddRoundKey, lalu InvMixColumns.
                    \item Putaran terakhir tanpa InvMixColumns.
                    \item Ditutup XOR dengan kunci awal.
                \end{itemize}
            \end{block}
        \end{column}
        \begin{column}{0.48\textwidth}
            \begin{exampleblock}{Transformasi invers}
                \begin{itemize}
                    \item InvSubBytes memakai S-box invers:
                          $\code{26} \to \code{23}$.
                    \item InvShiftRows menggeser baris ke kanan sejauh $r$.
                    \item InvMixColumns memakai matriks berkoefisien
                          \code{0E}, \code{0B}, \code{0D}, dan \code{09}.
                \end{itemize}
            \end{exampleblock}
            \begin{alertblock}{Simetri}
                Struktur enkripsi dan dekripsi sama; yang berubah adalah
                arah dan tabelnya.
            \end{alertblock}
        \end{column}
    \end{columns}
\end{frame}

% =====================================================================
\section{Contoh Kerja Transformasi}
% =====================================================================

\begin{frame}{Contoh: State Awal}
    \footnotesize
    State berikut dipakai pada tiga slide berikutnya. Seluruh nilai dalam
    heksadesimal dan diambil dari contoh Bab 8.
    \vspace{0.15cm}
    \[
        \text{state} =
        \begin{bmatrix}
            \code{23} & \code{A2} & \code{BC} & \code{4A} \\
            \code{D4} & \code{03} & \code{97} & \code{F3} \\
            \code{16} & \code{48} & \code{CD} & \code{50} \\
            \code{FF} & \code{DA} & \code{10} & \code{64}
        \end{bmatrix}
    \]
    \vspace{0.1cm}
    \begin{itemize}
        \item Baris 0 memuat \code{23}, \code{A2}, \code{BC}, \code{4A}.
        \item Baris 3 memuat \code{FF}, \code{DA}, \code{10}, \code{64}.
        \item Byte \code{23} berada di posisi $(0,0)$; byte \code{64} di
              posisi $(3,3)$.
    \end{itemize}
\end{frame}

\begin{frame}{Contoh: SubBytes}
    \footnotesize
    Setiap byte diganti nilai S-box pada baris dan kolom digitnya.
    \vspace{0.1cm}
    \[
        \begin{bmatrix}
            \code{23} & \code{A2} & \code{BC} & \code{4A} \\
            \code{D4} & \code{03} & \code{97} & \code{F3} \\
            \code{16} & \code{48} & \code{CD} & \code{50} \\
            \code{FF} & \code{DA} & \code{10} & \code{64}
        \end{bmatrix}
        \longrightarrow
        \begin{bmatrix}
            \code{26} & \code{3A} & \code{65} & \code{D6} \\
            \code{48} & \code{7B} & \code{88} & \code{0D} \\
            \code{47} & \code{52} & \code{BD} & \code{53} \\
            \code{16} & \code{57} & \code{CA} & \code{43}
        \end{bmatrix}
    \]
    \vspace{0.05cm}
    \begin{exampleblock}{Tiga pembacaan S-box}
        \scriptsize
        \code{23}: baris 2 kolom 3 $\to$ \code{26};\quad
        \code{A2}: baris A kolom 2 $\to$ \code{3A};\quad
        \code{FF}: baris F kolom F $\to$ \code{16}.
    \end{exampleblock}
    \vspace{0.02cm}
    Posisi setiap byte tetap; hanya nilainya yang berubah.
\end{frame}

\begin{frame}{Contoh: ShiftRows}
    \footnotesize
    Baris $r$ digeser siklik ke kiri sejauh $r$ byte.
    \vspace{0.1cm}
    \[
        \begin{bmatrix}
            \code{26} & \code{3A} & \code{65} & \code{D6} \\
            \code{7B} & \code{88} & \code{0D} & \code{48} \\
            \code{BD} & \code{53} & \code{47} & \code{52} \\
            \code{43} & \code{16} & \code{57} & \code{CA}
        \end{bmatrix}
        \longrightarrow
        \begin{bmatrix}
            \code{26} & \code{3A} & \code{65} & \code{D6} \\
            \code{88} & \code{0D} & \code{48} & \code{7B} \\
            \code{47} & \code{52} & \code{BD} & \code{53} \\
            \code{CA} & \code{43} & \code{16} & \code{57}
        \end{bmatrix}
    \]
    \vspace{0.05cm}
    \begin{itemize}
        \item Baris 0 tidak bergeser: \code{26}, \code{3A},
              \code{65}, \code{D6}.
        \item Baris 1 bergeser satu byte: \code{7B} pindah ke ujung kanan.
        \item Baris 3 bergeser tiga byte: \code{CA} pindah ke ujung kiri.
    \end{itemize}
    \begin{alertblock}{Hasil yang penting}
        \scriptsize
        Kolom 0 kini memuat \code{26}, \code{7B}, \code{BD}, \code{43}.
        SubBytes dan ShiftRows dapat ditukar urutannya, sehingga hasil ini
        setara dengan contoh ShiftRows pada buku.
    \end{alertblock}
\end{frame}

\begin{frame}{Contoh: MixColumns}
    \footnotesize
    Kolom 0 hasil ShiftRows adalah $[\code{26}\ \code{7B}\ \code{BD}\
    \code{43}]^{T}$.
    \vspace{0.05cm}
    \[
        \begin{bmatrix}
            \code{02} & \code{03} & \code{01} & \code{01} \\
            \code{01} & \code{02} & \code{03} & \code{01} \\
            \code{01} & \code{01} & \code{02} & \code{03} \\
            \code{03} & \code{01} & \code{01} & \code{02}
        \end{bmatrix}
        \begin{bmatrix}
            \code{26} \\ \code{7B} \\ \code{BD} \\ \code{43}
        \end{bmatrix}
        =
        \begin{bmatrix}
            \code{3F} \\ \code{4F} \\ \code{F9} \\ \code{2A}
        \end{bmatrix}
    \]
    \vspace{0.02cm}
    \begin{exampleblock}{Rincian baris pertama}
        \scriptsize
        $\code{02} \bullet \code{26} = \code{4C}$,\quad
        $\code{03} \bullet \code{7B} = \code{8D}$,\quad
        $\code{01} \bullet \code{BD} = \code{BD}$,\quad
        $\code{01} \bullet \code{43} = \code{43}$.\\[2pt]
        $\code{4C} \oplus \code{8D} = \code{C1}$,\quad
        $\code{C1} \oplus \code{BD} = \code{7C}$,\quad
        $\code{7C} \oplus \code{43} = \code{3F}$.
    \end{exampleblock}
    \vspace{0.02cm}
    Baris lain dihitung dengan pola yang sama dan menghasilkan
    \code{4F}, \code{F9}, serta \code{2A}.
\end{frame}

% =====================================================================
\section{Cipher Blok Lain}
% =====================================================================

\begin{frame}{GOST (Magma)}
    \footnotesize
    \begin{itemize}
        \item Blok 64 bit, kunci 256 bit, dan 32 putaran Jaringan Feistel.
        \item Standar Uni Soviet GOST 28147-89; standar baru memberinya
              nama Magma.
        \item Hanya ada delapan kunci putaran $K_1$ sampai $K_8$: dipakai
              berulang tiga kali, lalu urutannya dibalik pada delapan
              putaran terakhir.
        \item Fungsi $f$: jumlah modulo $2^{32}$ dengan kunci, substitusi
              delapan kotak-S $4 \times 4$ bit, lalu rotasi kiri 11 bit.
        \item Ciri pembeda: blok 64 bit dengan kunci 256 bit dan jadwal
              kunci yang sangat sederhana.
    \end{itemize}
    \vspace{0.05cm}
    \begin{alertblock}{Kunci panjang bukan jaminan}
        \footnotesize
        Serangan struktural menurunkan kompleksitas dari $2^{256}$ ke
        sekitar $2^{178}$. Untuk blok 64 bit, satu kunci juga dibatasi
        sekitar $2^{32}$ blok (kira-kira 32 GiB).
    \end{alertblock}
\end{frame}

\begin{frame}{RC5}
    \footnotesize
    \begin{columns}[T,onlytextwidth]
        \begin{column}{0.58\textwidth}
            \begin{itemize}
                \item Blok $2w$ bit: 32, 64, atau 128 bit untuk $w = 16,
                      32, 64$.
                \item Kunci 0 sampai 2040 bit; putaran 0 sampai 255.
                \item Jumlah kunci putaran $t = 2r + 2$; larik $S$
                      diinisialisasi dengan $P_w$ dan $Q_w$.
                \item Tiap putaran memakai XOR, rotasi sirkuler, dan
                      penjumlahan modulo $2^w$.
                \item Ciri pembeda: besar rotasi ditentukan oleh data
                      (nilai $B$ atau $A$), bukan konstanta tetap.
            \end{itemize}
        \end{column}
        \begin{column}{0.36\textwidth}
            \centering
            \bookfig[0.72\linewidth]{09}{putaran-rc5.png}
            \figsumber{Munir, \textit{IF4020 Kriptografi}, ITB.}
        \end{column}
    \end{columns}
\end{frame}

\begin{frame}{RC6}
    \footnotesize
    \begin{itemize}
        \item Blok 128 bit; kunci 128, 192, 256 sampai 2040 bit.
        \item Dikembangkan dari RC5 oleh Rivest, Robshaw, Sidney, dan Yin
              pada 1998; menjadi finalis AES.
        \item Memakai empat register $A$, $B$, $C$, dan $D$, masing-masing
              $w$ bit.
        \item Ciri pembeda: perkalian $f(x) = x(2x+1)$ ikut menentukan
              besar rotasi.
        \item Dapat dipandang sebagai dua RC5 yang bekerja paralel;
              penjadwalan kuncinya sama dengan RC5.
    \end{itemize}
    \vspace{0.1cm}
    \begin{exampleblock}{Yang bertambah dari RC5}
        \footnotesize
        Operasi dasarnya tetap XOR, penjumlahan, dan rotasi; RC6
        menambahkan perkalian dan satu pasang register.
    \end{exampleblock}
\end{frame}

\begin{frame}{Blowfish}
    \footnotesize
    \begin{itemize}
        \item Blok 64 bit; kunci 32 sampai 448 bit; 16 putaran Jaringan
              Feistel.
        \item Dirancang Bruce Schneier pada 1993 dan memakai empat kotak-S.
        \item Fungsi $F$ memecah 32 bit menjadi empat byte $a,b,c,d$:
              \[
                  F(x) = \bigl((S_1[a] \boxplus S_2[b]) \oplus S_3[c]\bigr)
                  \boxplus S_4[d].
              \]
        \item Ciri pembeda: empat kotak-S dan penyaringan keluaran dengan
              $K_{17}$ serta $K_{18}$ setelah 16 putaran.
    \end{itemize}
    \vspace{0.05cm}
    \begin{alertblock}{Penyaringan keluaran}
        \footnotesize
        Pertukaran terakhir pada Jaringan Feistel dibatalkan lebih dahulu;
        enkripsi kemudian ditutup dengan dua kunci tambahan itu.
    \end{alertblock}
\end{frame}

\begin{frame}{Serpent}
    \footnotesize
    \begin{itemize}
        \item Blok 128 bit; kunci 128, 192, atau 256 bit; 32 putaran.
        \item Dirancang Anderson, Biham, dan Knudsen; finalis AES.
        \item Setiap putaran melakukan XOR dengan subkunci, substitusi
              kotak-S, lalu transformasi linear.
        \item Memakai delapan kotak-S, permutasi awal, dan permutasi akhir.
        \item Ciri pembeda: 32 putaran yang dibingkai permutasi awal dan
              permutasi akhir.
    \end{itemize}
    \vspace{0.1cm}
    \begin{exampleblock}{Penilaian Bab 9}
        \footnotesize
        Tabel perbandingan cipher simetris menempatkan Serpent pada
        kategori sangat kuat.
    \end{exampleblock}
\end{frame}

\begin{frame}{Twofish}
    \footnotesize
    \begin{columns}[T,onlytextwidth]
        \begin{column}{0.56\textwidth}
            \begin{itemize}
                \item Blok 128 bit; kunci sampai 256 bit; 16 putaran.
                \item Finalis AES, dirancang Bruce Schneier, dan
                      dipandang sebagai kelanjutan Blowfish.
                \item Memakai empat kotak-S $8 \times 8$ yang bergantung
                      pada kunci.
                \item Ciri pembeda: dua fungsi $g$, matriks MDS, dan
                      transformasi PHT dalam setiap putaran.
            \end{itemize}
        \end{column}
        \begin{column}{0.38\textwidth}
            \centering
            \bookfig[0.70\linewidth]{09}{putaran-twofish.png}
            \figsumber{Munir, \textit{IF4020 Kriptografi}, ITB.}
        \end{column}
    \end{columns}
\end{frame}

\begin{frame}{MARS}
    \footnotesize
    \begin{itemize}
        \item Blok 128 bit; kunci 128 sampai 448 bit dengan kenaikan
              32 bit; 32 putaran.
        \item Finalis AES yang dirancang IBM.
        \item Bukan cipher substitusi--permutasi biasa, melainkan
              Jaringan Feistel tak seimbang.
        \item Ciri pembeda: 32 putaran terbagi menjadi 8 putaran
              \textit{forward mixing}, 16 putaran inti, dan 8 putaran
              \textit{backward mixing}.
        \item Fase pencampuran memakai kotak-S $8 \times 32$, pergeseran
              8 bit, penjumlahan, dan XOR.
    \end{itemize}
    \vspace{0.1cm}
    \begin{alertblock}{Penyaringan}
        \footnotesize
        Penyaringan memakai $K_0$ sampai $K_3$ serta $K_{36}$ sampai
        $K_{39}$.
    \end{alertblock}
\end{frame}

% =====================================================================
\section{Perbandingan dan Penutup}
% =====================================================================

\begin{frame}{Perbandingan Cipher Blok Pertemuan Ini}
    \scriptsize
    \begin{tabularx}{\linewidth}{@{}l c c c >{\raggedright\arraybackslash}X@{}}
        \toprule
        \textbf{Cipher} & \textbf{Blok} & \textbf{Kunci} & \textbf{Putaran} &
        \textbf{Ciri pembeda} \\
        \midrule
        AES & 128 & 128/192/256 & 10/12/14 & Substitusi--permutasi pada byte \\
        GOST & 64 & 256 & 32 & Feistel, 8 kotak-S $4 \times 4$ \\
        RC5 & $2w$ & 0--2040 & 0--255 & Rotasi bergantung data \\
        RC6 & 128 & 128--2040 & $r$ & Perkalian $x(2x+1)$ \\
        Blowfish & 64 & 32--448 & 16 & Kotak-S bergantung kunci \\
        Serpent & 128 & 128/192/256 & 32 & Delapan kotak-S, 32 putaran \\
        Twofish & 128 & s.d. 256 & 16 & MDS dan PHT \\
        MARS & 128 & 128--448 & 32 & Feistel tak seimbang 8+16+8 \\
        \bottomrule
    \end{tabularx}
    \vspace{0.15cm}
    \begin{alertblock}{Dua ukuran yang menentukan pemakaian}
        \footnotesize
        Blok 64 bit membatasi data per kunci sekitar $2^{32}$ blok;
        blok 128 bit memberi ruang yang jauh lebih besar.
    \end{alertblock}
\end{frame}

\begin{frame}{Rangkuman (1/2)}
    \footnotesize
    \begin{itemize}
        \item AES adalah Rijndael karya Daemen dan Rijmen; bloknya 128 bit,
              kuncinya 128/192/256 bit, dan putarannya 10/12/14.
        \item Data selama enkripsi disimpan pada state $4 \times 4$ byte;
              masukan disalin kolom demi kolom.
        \item SubBytes mengganti byte lewat S-box, ShiftRows menggeser
              baris, MixColumns mengalikan kolom di $\text{GF}(2^8)$, dan
              AddRoundKey melakukan XOR dengan kunci putaran.
        \item Alur AES-128: AddRoundKey awal, sembilan putaran penuh, lalu
              putaran terakhir tanpa MixColumns.
        \item Key expansion memakai RotWord, SubWord, dan konstanta
              $Rcon$ untuk membangkitkan sebelas kunci putaran.
    \end{itemize}
\end{frame}

\begin{frame}{Rangkuman (2/2)}
    \footnotesize
    \begin{itemize}
        \item Ruang kunci AES-128 berisi $2^{128} \approx 3{,}4 \times
              10^{38}$ kemungkinan, sehingga penelusuran menyeluruh tidak
              praktis.
        \item Contoh buku: state \code{23}\,\code{A2}\,\code{BC}\,\code{4A}
              dan seterusnya, setelah SubBytes dan ShiftRows menghasilkan
              kolom \code{26}\,\code{7B}\,\code{BD}\,\code{43}, lalu
              MixColumns memberi \code{3F}\,\code{4F}\,\code{F9}\,\code{2A}.
        \item GOST: blok 64 bit, kunci 256 bit, 32 putaran Feistel.
        \item RC5 dan RC6 menonjol karena parameter dan rotasi yang
              lentur; Blowfish dan Twofish memakai kotak-S bergantung
              kunci.
        \item Serpent menekankan banyak putaran, sedangkan MARS memakai
              Feistel tak seimbang tiga fase.
    \end{itemize}
\end{frame}

\begin{frame}{Penutup}
    \begin{block}{Yang harus diingat}
        \footnotesize
        \begin{itemize}
            \item AES bukan jaringan Feistel dan bekerja pada byte.
            \item Empat transformasi putaran memiliki peran berbeda:
                  substitusi, pergeseran, penyebaran, dan penyisipan kunci.
            \item Panjang kunci bukan satu-satunya ukuran keamanan;
                  struktur dan batas \textit{rekeying} ikut menentukan.
            \item Cipher blok lain dibedakan oleh ukuran blok, ukuran
                  kunci, jumlah putaran, dan struktur putarannya.
        \end{itemize}
    \end{block}
    \vspace{0.2cm}
    \begin{exampleblock}{Pertemuan berikutnya}
        \textbf{Minggu 11 --- Algoritma RSA}: properti, pembangkitan kunci,
        enkripsi dan dekripsi, keamanan, kelemahan, serta
        \textit{man-in-the-middle}.
    \end{exampleblock}
\end{frame}

\end{document}
